% 第 23 讲 第二不等式的证明（源：src/lecture23.txt；已删去抽出时残留的页眉「23. PROOF OF THE SECOND INEQUALITY 103/105/107」）
\section*{第 23 讲\quad 第二不等式的证明}

本讲的目标是证明「第二不等式」：对整体域的一切扩张 $E/F$，有 $H^1(G,C_E)=0$，其中 $C_E:=A_E^\times/E^\times$ 是 $E$ 的「伊代尔类群」。我们的主要情形是 $E/F$ 为 $p$ 阶循环扩张，且对某个本原 $p$ 次单位根 $\zeta_p$ 有 $\zeta_p\in F$；我们将在本讲末尾归约到此情形（注意 $\operatorname{char}(F)=0$，因为我们假设 $F$ 是数域）。在此情形下，只需证明
\[
\#\widehat H^0(C_E)=\#(C_F/NC_E)=\#\bigl(A_F^\times/F^\times\cdot N(A_E^\times)\bigr)\le p.
\]
事实上，由「第一不等式」我们知道
\[
\frac{\#\widehat H^0(C_E)}{\#\widehat H^1(C_E)}=p,
\]
于是 $p\cdot\#\widehat H^1(C_E)=\#\widehat H^0(C_E)\le p$ 蕴含 $\#H^1(C_E)=1$，正如所愿。我们的方法将是「试错」式的——即先尝试某种做法，它还不够好，然后再加以修正。

取定 $F$ 的一个有限素位集 $S$，使得
\begin{enumerate}[label=(\arabic*),leftmargin=2.4em]
\item 若 $v\mid\infty$，则 $v\in S$；
\item 若 $v\mid p$，则 $v\in S$；
\item $A_F^\times=F^\times\cdot A_{F,S}^\times$，其中回忆
\[
A_{F,S}^\times:=\prod_{v\in S}F_v^\times\times\prod_{v\notin S}\mathcal{O}_{F_v}^\times,
\]
且由引理 20.12 这是可行的；
\item $E=F(\sqrt[p]{u})$，其中 $u\in\mathcal{O}_{F,S}^\times:=F^\times\cap A_{F,S}^\times$ 是 $F$ 的「$S$-单位」。
\end{enumerate}
最后一条由 Kummer 理论可行。

注意最后这个条件蕴含 $E$ 在 $S$ 外非分歧：因为 $u$ 在任何 $v\notin S$ 处是整元素，且由假设所有整除 $p$ 的素位都在 $S$ 中，$p$ 与 $F_v$ 的剩余域的阶互素，故 $F_v(\sqrt[p]{u})/F_v$ 是非分歧扩张。

一个重要的断言（稍后将以略为精细的形式证明）如下：

\textbf{断言 23.1.} $u\in\mathcal{O}_{F,S}^\times$ 是 $p$ 次幂，当且仅当它在 $F_v^\times$ 中的像对每个 $v\in S$ 都是 $p$ 次幂。

令
\[
\Gamma:=\prod_{v\in S}(F_v^\times)^p\times\prod_{v\notin S}\mathcal{O}_{F_v}^\times\subseteq A_{F,S}^\times.
\]
则我们有下述断言：

\textbf{断言 23.2.} $\mathcal{O}_{F,S}^\times\cap\Gamma=(\mathcal{O}_{F,S}^\times)^p$。

\textbf{证明.} 这由前一断言平凡地推出。\hfill$\square$

\textbf{断言 23.3.} $\Gamma\subseteq N(A_E^\times)$。

\textbf{证明.} 扩张 $E/F$ 在每个 $v\notin S$ 处非分歧，故因子 $\prod_{v\notin S}\mathcal{O}_{F_v}^\times\subseteq N(A_E^\times)$。由于 $p$ 杀掉了 $\widehat H^0(E_w^\times)$（对某个选定的 $w\mid v$），可知因子 $\prod_{v\in S}(F_v^\times)^p\subseteq N(A_E^\times)$ 亦然。\hfill$\square$

于是
\[
\#\bigl(A_F^\times/F^\times\cdot N(A_E^\times)\bigr)\le\#\bigl(A_F^\times/F^\times\cdot\Gamma\bigr),
\]
且有短正合列
\[
1\to\mathcal{O}_{F,S}^\times/(\mathcal{O}_{F,S}^\times\cap\Gamma)\to A_{F,S}^\times/\Gamma\to A_F^\times/(F^\times\cdot\Gamma)\to 1.
\]
事实上，第三个映射由上面 $S$ 的性质 (3) 满射，第二个映射单射因为 $\mathcal{O}_{F,S}^\times\subseteq F^\times$，而在 $A_{F,S}^\times/\Gamma$ 处的正合性按定义成立。于是
\[
\#\bigl(A_F^\times/F^\times\cdot\Gamma\bigr)=\frac{\#(A_{F,S}^\times/\Gamma)}{\#(\mathcal{O}_{F,S}^\times/\mathcal{O}_{F,S}^\times\cap\Gamma)},
\]
剩下的就是计算此式的分子与分母。我们有
\[
A_{F,S}^\times/\Gamma=\prod_{v\in S}F_v^\times/(F_v^\times)^p,
\]
% [原文此处作 "∏_{p∈S}"，而通项用 v；按上下文作 v∈S，见报告 §3]
并由 (6.3) 回忆
\[
\#\bigl(F_v^\times/(F_v^\times)^p\bigr)=p\cdot\frac{\#\mu_p(F_v)}{|p|_v}=\frac{p^2}{|p|_v},
\]
因为由假设 $\zeta_p\in F$。于是
\[
\prod_{v\in S}\frac{p^2}{|p|_v}=p^{2\#S}
\]
由乘积公式得到，因为由假设 $v\notin S$ 时 $|p|_v=1$。现在想计算
\[
\#\bigl(\mathcal{O}_{F,S}^\times/\mathcal{O}_{F,S}^\times\cap\Gamma\bigr)=\#\bigl(\mathcal{O}_{F,S}^\times/(\mathcal{O}_{F,S}^\times)^p\bigr).
\]
回忆由 $S$-单位定理，
\[
\mathcal{O}_{F,S}^\times\cong\Z^{\#S-1}\times(\mathcal{O}_{F,S}^\times)_{\mathrm{tors}}.
\]
后者循环，且其阶被 $p$ 整除，故
\[
\#\bigl(\mathcal{O}_{F,S}^\times/(\mathcal{O}_{F,S}^\times)^p\bigr)=p^{\#S-1}\cdot p=p^{\#S}.
\]
合并这两个结果，得到
\[
\#\widehat H^0(C_E)\le\frac{p^{2\#S}}{p^{\#S}}=p^{\#S},
\]
遗憾的是这还不够好。

下面是我们改进此结果的办法：

\textbf{断言 23.4.} 给定这样的集合 $S\subseteq M_F$，存在集合 $T\subseteq M_F$，使得
\begin{enumerate}[label=(\arabic*),leftmargin=2.4em]
\item $\#T=\#S-1$；
\item $S\cap T=\varnothing$；
\item 每个 $v\in T$ 在 $E$ 中分裂，即对一切 $w\mid v$ 有 $E_w=F_v$；
\item 任意 $u\in\mathcal{O}_{F,S\cup T}^\times$ 是 $p$ 次幂，当且仅当 $u\in F_v^\times$ 对一切 $v\in S$ 都是 $p$ 次幂。
\end{enumerate}

注意这里与先前情形的关键差别：在性质 (4) 中，我们不要求 $u\in(F_v^\times)^p$ 对一切 $v\in S\cup T$ 成立，只要求对一切 $v\in S$ 成立。给定这样的 $T$，我们重新定义 $\Gamma$ 为
\[
\Gamma:=\prod_{v\in S}(F_v^\times)^p\times\prod_{v\in T}F_v^\times\times\prod_{v\notin S\cup T}\mathcal{O}_{F_v}^\times.
\]

\textbf{断言 23.5.} $\Gamma\subseteq N(A_E^\times)$。

\textbf{证明.} 性质 (3) 给出关于第二个因子的断言；第一个与第三个因子如前。\hfill$\square$

用以 $A_{F,S\cup T}^\times$ 代替 $A_{F,S}^\times$ 的方式重做我们的计算，由性质 (4) 如前得到
\[
\#(A_{F,S\cup T}^\times/\Gamma)=p^{2\#S},
\]
并且同样如前，
\[
\#\bigl(\mathcal{O}_{F,S\cup T}^\times/(\mathcal{O}_{F,S\cup T}^\times\cap\Gamma)\bigr)=p^{\#(S\cup T)}=p^{2\#S-1},
\]
故二者之商为 $p$，正如所愿！于是只需证明上述断言。

\textbf{断言 23.6.} 对整体域的任意阿贝尔扩张 $F'/F$，$v\notin S$ 的 Frobenius 元素生成 $\Gal(F'/F)$。

我们希望纯粹代数地证明这一点，而不用切博塔廖夫密度定理（后者无论如何给出的是略为不同的陈述）。

\textbf{证明.} 设 $H$ 是由一切 $v\notin S$ 的 Frobenius 生成的子群，$F'':=(F')^H$ 为不动域。我们希望证明 $F''=F$。注意对一切 $v\notin S$，$\operatorname{Frob}_v$ 在 $\Gal(F'/F)/H=\Gal(F''/F)$ 中平凡，故每个 $v\notin S$ 都在 $F''/F$ 中分裂（因为由假设它们非分歧）。于是对一切 $w\mid v$ 与 $v\notin S$ 有 $F''_w=F_v$，而我们断言这不可能。

可设 $F''/F$ 是 $n$ 次循环扩张（必要时以更小的扩张代替之）。由第一不等式，$\chi(C_{F''})=n$，这给出
\[
\#\bigl(A_F^\times/N(A_{F''}^\times)\cdot F^\times\bigr)=\#\widehat H^0(C_{F''})\ge n.
\]
但因此扩张对一切 $v\notin S$ 分裂，平凡地有 $N((F''_v)^\times)=F_v^\times$，从而 $\prod\coprod_{v\notin S}F_v^\times\subseteq N(A_{F''}^\times)$，其中这是限制直积。强逼近随后给出 $F^\times\cdot\prod\coprod_{v\notin S}F_v^\times$ 在 $A_F^\times$ 中稠密，而它又是开的，故除非 $n=1$ 便得到矛盾，正如所愿。\hfill$\square$

我们希望对此断言应用于 $F':=F(\{\sqrt[p]{u}:u\in\mathcal{O}_{F,S}^\times\})$。先有一个断言：

\textbf{断言 23.7.} 对上述 $F'$，有 $\Gal(F'/F)=(\Z/p\Z)^{\#S}$。

\textbf{证明.} 这实质上是 Kummer 理论，因为 $\mathcal{O}_{F,S}^\times/(\mathcal{O}_{F,S}^\times)^p\subseteq F^\times/(F^\times)^p$。我们知道 $F$ 的一切指数 $p$ 的扩张都通过添加 $F^\times$ 中元素的 $p$ 次根得到。其 Galois 群必是 $\Z/p\Z$ 的若干个拷贝之积，但其中某些子群可能重合——迭代应用 Kummer 理论便给出该陈述。\hfill$\square$

现在我们有 $F'/E/F$，因为假设 $E/F$ 是通过添加某个 $S$-单位的 $p$ 次根得到的。选取 $E$ 的素位 $w_1,\dots,w_{\#S-1}$，它们不整除 $S$ 中任何素位，且其 Frobenius 给出 $\Gal(F'/E)\cong(\Z/p\Z)^{\#S-1}$ 的一组基；由断言 23.6 的论证这是可能的。然后令 $T:=\{v_1,\dots,v_{\#S-1}\}$ 为各 $w_i$ 到 $F$ 的限制。

\textbf{断言 23.8.} 每个 $v\in T$ 在 $E$ 中分裂。

\textbf{证明.} 由于 $\operatorname{Frob}_v\in\Gal(F'/E)$，它在 $E$ 上平凡作用，故对任意 $w\mid v$，$\Gal(E_w/F_v)$ 平凡，正如所愿。\hfill$\square$

这建立了 $T$ 的条件 (3)；还需证明条件 (4)，因为条件 (1) 与 (2) 在 $T$ 的构造中已隐含。

\textbf{断言 23.9.} 元素 $x\in\mathcal{O}_{F,S\cup T}^\times$ 是 $p$ 次幂，当且仅当对每个 $v\in S$ 有 $x\in(F_v^\times)^p$。

\textbf{证明.} \emph{第一步.} 我们断言
\[
\mathcal{O}_{F,S}^\times\cap(E^\times)^p=\{x\in\mathcal{O}_{F,S}^\times:x\in(F_v^\times)^p\ \text{对一切 } v\in T\}.
\]
正向包含平凡，因为由前一断言 $(F_v^\times)^p=(E_w^\times)^p$。对反向，注意对任意 $x\in\mathcal{O}_{F,S}^\times$，有扩张 $F'/E(\sqrt[p]{x})/E$。若对每个 $w\mid v$ 与 $v\in T$ 有 $x\in(E_w^\times)^p$，则此扩张在 $w$ 处分裂，故 $\operatorname{Frob}_w$ 在 $E(\sqrt[p]{x})$ 上平凡作用；于是 $\Gal(F'/E)$ 在 $E(\sqrt[p]{x})$ 上平凡作用，因为它由这些 Frobenius 生成；故 $E(\sqrt[p]{x})=E$ 且 $x\in(E^\times)^p$，正如所愿。

\emph{第二步.} 现在我们断言典范映射
\[
\mathcal{O}_{F,S}^\times\xrightarrow{\ \varphi\ }\prod_{v\in T}\mathcal{O}_{F_v}^\times/(\mathcal{O}_{F_v}^\times)^p
\]
满射。正是这一步真正建立了对 $T$ 的大小的限制，第二不等式由此完美地推出。我们将通过计算两端的阶来进行。由第一步，
\[
\operatorname{Ker}(\varphi)=\{x\in\mathcal{O}_{F,S}^\times:x\in(E^\times)^p\}.
\]
则 $\mathcal{O}_{F,S}^\times/\operatorname{Ker}(\varphi)$ 的阶为 $p^{\#S-1}$。事实上，我们先前算出 $\mathcal{O}_{F,S}^\times/(\mathcal{O}_{F,S}^\times)^p$ 的阶为 $p^{\#S}$，且由于
\[
(\mathcal{O}_{F,S}^\times)^p=\{x\in\mathcal{O}_{F,S}^\times:x\in(F^\times)^p\}
\]
而 $E/F$ 是添加某个 $S$-单位的 $p$ 次根得到的 $p$ 次扩张，可知 $[\operatorname{Ker}(\varphi):(\mathcal{O}_{F,S}^\times)^p]=p$。现在，用我们先前的公式关于 $\mathcal{O}_{F_v}^\times$（而非 $F_v^\times$）的版本，右端的阶为
\[
\prod_{v\in T}\frac{\#\mu_p(F_v)}{|p|_v}=p^{\#T}=p^{\#S-1},
\]
故该映射确实满射。

\emph{第三步.} 现在希望建立此断言：若对一切 $v\in S$ 有 $x\in(F_v^\times)^p$，则 $x\in(\mathcal{O}_{F,S\cup T}^\times)^p$（反向平凡）。我们希望证明 $F(\sqrt[p]{x})=F$。令
\[
\Gamma:=\prod_{v\in S}F_v^\times\times\prod_{v\in T}(\mathcal{O}_{F_v}^\times)^p\times\prod_{v\notin S\cup T}\mathcal{O}_{F_v}^\times\subseteq A_{F,S}^\times,
\]
这里又是一个与先前不同的 $\Gamma$。则事实上
\[
\Gamma\subseteq N(A_{F(\sqrt[p]{x})}^\times)\subseteq A_F^\times,
\]
其中第三个因子是因为 $F(\sqrt[p]{x})/F$ 在 $S\cup T$ 外非分歧，第二个因子是因为 $[F(\sqrt[p]{x}):F]\le p$，第一个因子是因为由假设该扩张在 $S$ 的一切素位处分裂。现在我们要证明 $F^\times\cdot\Gamma=A_F^\times$，因为第一不等式随后便如断言 23.6 那样给出结果。由第二步，有
\[
\mathcal{O}_{F,S}^\times\twoheadrightarrow\prod_{v\in T}\mathcal{O}_{F_v}^\times/(\mathcal{O}_{F_v}^\times)^p=A_{F,S}^\times/\Gamma,
\]
故 $\mathcal{O}_{F,S}^\times\cdot\Gamma=A_{F,S}^\times$。这意味着
\[
F^\times\cdot\Gamma=F^\times\cdot A_{F,S}^\times=A_F^\times
\]
由关于 $S$ 的假设。\hfill$\square$

现在希望由已证的特殊情形推出第二不等式的一般情形。第一步如下：

\textbf{断言 23.10.} 若第二不等式对任意基域含 $p$ 次单位根的 $p$ 阶循环扩张成立，则它对任意 $p$ 阶循环扩张成立。

\textbf{证明.} 设 $E/F$ 是整体域的 $p$ 次循环扩张。回忆 $E/F$ 的第二不等式等价于典范单射
\[
\Br(F/E)\hookrightarrow\bigoplus_{v\in M_F}\Br(F_v)
\]
的存在性。事实上，有短正合列
\[
0\to E^\times\to A_E^\times\to C_E\to 0,
\]
而上同调长正合列于是给出
\[
\underbrace{H^1(G,A_E^\times)}_{\oplus H^1(E_w^\times)=0}\to H^1(G,C_E)\to\Br(F/E)\to\bigoplus_v\Br(F_v/E_w)\subseteq\bigoplus_v\Br(F_v)
\]
（对某个选定的 $w\mid v$），其中第一个等号由希尔伯特定理 90 得到。为证明 $H^1(G,C_E)$ 的消没，只需证明最后的映射单射。现在，域扩张
\[
\begin{array}{ccccc}
& & E(\zeta_p) & &\\
& \nearrow & & \nwarrow &\\
E & & & & F(\zeta_p)\\
& \nwarrow & & \nearrow &\\
& & F & &
\end{array}
\]
诱导出交换图
\[
\begin{array}{ccc}
\Br(F/E) & \xrightarrow{\ \alpha\ } & \bigoplus_v\Br(F_v/E_w)\\[4pt]
\big\downarrow{\scriptstyle\gamma} & & \big\downarrow{\scriptstyle\delta}\\[4pt]
\Br(F(\zeta_p)/E(\zeta_p)) & \xrightarrow{\ \beta\ } & \bigoplus_v\Br(F(\zeta_p)_w),\\[4pt]
\big\downarrow & & \\[4pt]
\Br(F/E) & &
\end{array}
\]
% [按原 PDF 第 106 页复原：原图在最左端用花括号把左上角的 $\Br(F/E)$ 与左下角的 $\Br(F/E)$ 括在一起，并标以 $\times[F(\zeta_p):F]$，即左列两个竖直映射（限制与膨胀）的复合等于乘以 $[F(\zeta_p):F]$]
其中最左的两个映射分别是上同调上的限制映射与膨胀映射，其依据是布劳尔群的上同调解释（见习题集 7 第 2 题）。此外，左列复合 $\Br(F/E)\to\Br(F(\zeta_p)/E(\zeta_p))\to\Br(F/E)$ 是乘以 $[F(\zeta_p):F]$，它在 $\Br(F/E)$ 上单射，因为后者是 $p$-挠的（由习题 2(c)），且 $[F(\zeta_p):F]\mid(p-1)$。于是 $\gamma$ 也单射。由于第二不等式对 $E(\zeta_p)/F(\zeta_p)$ 由假设成立，$\beta$ 单射，从而 $\alpha$ 也单射。
% [原文作 "Since the second equality holds for E(ζ_p)/F(ζ_p)"，此处应为 second inequality，见报告 §3]
\hfill$\square$

\textbf{断言 23.11.} 若第二不等式对任意数域的 $p$ 阶循环扩张成立，则它对任意扩张成立。

\textbf{证明.} 我们希望证明 $H^1(G,C_E)=0$。如同有限群的任何 Tate 上同调群，有单射
\[
H^1(G,C_E)\hookrightarrow\bigoplus_p H^1(G_p,C_E),
\]
其中 $G_p$ 是 $G$ 的 $p$-Sylow 子群。于是可设 $G$ 是 $p$-群。由于每个 $p$-群 $G$ 都含有同构于 $\Z/p\Z$ 的正规子群 $H$，可设我们有域扩张 $E_2/E_1/F$，其中 $\Gal(E_2/E_1)\cong H$，$\Gal(E_1/F)\cong G/H$。可设定理对 $H$ 作用在 $E_2$ 上以及 $G/H$ 作用在 $E_1$ 上成立，于是只需重复前一断言证明中关于布劳尔群的单射性的那类论证。

首先我们断言 $C_{E_2}^H=C_{E_1}$。事实上，有短正合列
\[
0\to E_2^\times\to A_{E_2}^\times\to C_{E_2}\to 0,
\]
而上同调长正合列于是给出
\[
0\to\underbrace{H^0(H,E_2^\times)}_{E_1^\times}\to\underbrace{H^0(H,A_{E_2}^\times)}_{A_{E_1}^\times}\to\underbrace{H^0(H,C_{E_2})}_{C_{E_2}^H}\to\underbrace{H^1(H,E_2^\times)}_{0}
\]
由希尔伯特定理 90。注意 $A_{E_2}^{\times,H}=A_{E_1}^\times$，因为取有限群的不变子模与阿代尔定义中的正向极限和直积交换。于是有
\[
\operatorname{hKer}\Bigl(C_{E_2}^{hG}=(C_{E_2}^{hH})^{hG/H}\to(\tau^{\ge2}C_{E_2}^{hH})^{hG/H}\Bigr)\cong(\tau^{\le0}C_{E_2}^{hH})^{hG/H}=(C_{E_1})^{hG/H},
\]
其中第一个等号由习题集 6 第 3 题，映射由截断的定义而来，拟同构是因为 $H^1(H,C_{E_2})$ 由假设消没，而最后这个表示式不过是 $C_{E_2}$ 的朴素 $H$-不变子模，因为截断杀掉了次数大于 $0$ 的一切上同调，故前一断言给出等式。上同调长正合列随后给出
\[
\underbrace{H^1\bigl((C_{E_1})^{hG/H}\bigr)}_{H^1(G/H,C_{E_1})=0}\to\underbrace{H^1\bigl((C_{E_2})^{hG}\bigr)}_{H^1(G,C_{E_2})}\to\underbrace{H^1\bigl((\tau^{\ge2}C_{E_2}^{hH})^{hG/H}\bigr)}_{0},
\]
因为最右的复形处于次数至少 $2$。于是 $H^1(G,C_{E_2})=0$，正如所愿。\hfill$\square$
